\section{Discussion: Policy Design}\label{sec:discussion}

The model puts a number on a trade-off that procurement authorities in Brazil are already making implicitly. Law 14{,}133/2021---the federal procurement framework adopted in 2021 and now being rolled out across the 5{,}570 Brazilian municipalities---contemplates both full SME set-asides and price preferences as tools for small-firm protection. The choice between the two instruments is currently decided at the administrative level, with the set-aside being the default for low-value procurement and the preference left to discretionary use. The estimates here indicate that the defaults should invert.

\subsection{The welfare weight implicit in the current regime}

Consider a social planner aggregating producer and consumer surplus with a weight $w^{\text{SME}}$ on SME producer surplus. The planner is indifferent between the full set-aside (V0) and the 10 percent price preference (V3) when the welfare-loss differential equals the SME-surplus differential, times the weight:
\begin{equation}
w^{\text{SME}}_\star
\;=\; \frac{\mathrm{Loss}(V_0) - \mathrm{Loss}(V_3)}
           {\mathrm{SMESurplus}(V_0) - \mathrm{SMESurplus}(V_3)}.
\end{equation}
With the point estimates of Section~\ref{sec:welfare} and the SME-surplus implicit in the V0--V3 comparison ($\mathrm{Loss}(V_0) - \mathrm{Loss}(V_3) \approx 0.22$ in non-pharma and $0.30$ in pharma; $\mathrm{SMESurplus}(V_0) - \mathrm{SMESurplus}(V_3) \approx 0.09$ and $0.10$), the indifference weight is
\[
w^{\text{SME}}_\star \approx 2.4 \text{ (non-pharma)}, \qquad 3.0 \text{ (pharma)}.
\]
For the full set-aside to be welfare-optimal under these estimates, the planner would have to value one real of SME surplus at between 2.4 and 3.0 reais of general welfare.

This is a large number. Brazilian cash-transfer programs (Bolsa Fam\'{\i}lia and Aux\'{\i}lio Brasil) have been defended under implicit weights in the 1.2--1.5 range---the shadow price of a real transferred to the lower tail of the income distribution, relative to a real in general tax revenue. The SMEs receiving the set-aside benefit are medium-sized formal-sector firms, not individuals in the lower tail; the standard welfare argument for concentrating fiscal resources on them is weaker than for cash transfers, not stronger. A weight of 2.4--3.0 sits above the weight implicit in Brazilian social policy for the population groups with the most established claim to such support. The set-aside is, on these numbers, Pareto-inefficient for any plausible welfare aggregator. I read the gap from 2.4--3.0 to the 1.2--1.5 benchmark as prima facie evidence that Brazil's procurement-protection regime is under-delivering relative to the 10 percent price preference.

\subsection{Implementation considerations}

Two features qualify the main result. First, V3 is computed in a Vickrey-equivalent framework. In a first-price sealed-bid auction with the same preference rule, bidder equilibrium behavior would adjust---SMEs would bid more aggressively anticipating the preference, non-SMEs would bid less aggressively anticipating the lost share---in ways characterized by \citet{maskinriley2000}. The adjustment reduces the realized gain to SMEs relative to the Vickrey benchmark, and increases the realized welfare saving to the government. The direction strengthens the Pareto-improvement claim; the magnitude depends on the strategic adjustment, which this paper does not solve.

Second, a preference rule is procedurally more complex to implement than a set-aside. The buyer must score bids twice (actual, and actual scaled by $1 - 0.10$ for SMEs), and the winner-selection rule must be disclosed and contestable. In the US and in the SBA framework the preference rule has been operational since the 1970s without significant administrative friction. The same is true for the \citet{denes1997} subcontracting mandate, a variant that achieves similar SME protection via a different mechanism. Brazilian state procurement has already experimented with preference rules in several narrow product categories---the scale-up to Group-65 medical supplies, and beyond, is a procedural rather than a conceptual move.

\subsection{Non-pharma versus pharma}

The welfare cost of the set-aside is 60 percent larger in pharma than in non-pharma---28 percent of procurement value versus 46 percent. Two forces contribute.

Auction-level heterogeneity is larger in pharma (intraclass correlation 0.62 versus 0.47). The UH-clean cost distribution has a heavier right tail in pharma, which amplifies the allocative wedge: when the rule forces the winner to come from the SME distribution only, the pharma winner's excess production cost is larger in expectation.

The SME pool in pharma is thinner in absolute terms ($N^{\text{SME}}_{\text{Pre}} = 0.44$ in pharma versus $0.86$ in non-pharma), so the policy change moves $N^{\text{SME}}$ from a much smaller base. The equilibrium price under SME-only, $p_{S_2}$, is pulled further from the open-regime $p_{S_1}$ when the SME pool is thin, because the second-order statistic is drawn from a thinner distribution. This is an intensive-margin effect amplified by thin SME baselines.

The case for relaxing the set-aside in pharma is therefore strongest. This matches the reduced-form reading in the v1 predecessor paper---that the CMED-regulated pharma items bore roughly twice the welfare distortion of the non-pharma medical subsegment---and strengthens it through the structural framework. Two Brazilian regulatory regimes (the CMED price ceiling and the SME set-aside) interact: the set-aside raises prices within the residual slack below the CMED cap, and the allocative cost is larger than it would be under open competition at the cap. The price ceiling does not shield the welfare account from the set-aside.

\subsection{Extensions the paper does not attempt}

The model treats procurement quantity $Q$ as inelastic and the set-aside as static. A full welfare account would price two more margins: the quantity response of the buyer to higher prices (typical in procurement with a fixed budget), and the dynamic capacity-accumulation effects for SMEs protected by the rule \citep{szerman2023,cabral2017}. Both extensions would typically strengthen the case against the full set-aside---quantity adjustment introduces additional deadweight beyond the static channels, and the dynamic SME capacity effects documented in Brazilian data are empirically small and short-lived. A full dynamic-general-equilibrium treatment of the SME-only rule is a paper of its own.
